\documentclass[tikz,border=6pt]{standalone}
\usepackage{peabra}
\begin{document}
\begin{tikzpicture}[x=1cm,y=1cm, l/.style={doboz=34mm, font=\scriptsize, minimum height=20mm, inner sep=2pt}]
  \foreach \i/\t/\d in {1/{Az elemzés céljának\\meghatározása}/{veszteségek (L),\\rendszerszintű veszélyek (H),\\biztonsági korlátok (SC)},
      2/{Az irányítási struktúra\\modellezése}/{irányítók, akciók,\\visszacsatolások,\\felelősségek},
      3/{Nem biztonságos irányítási\\akciók azonosítása}/{négy típus:\\nem ad, ad, rossz időben,\\rossz ideig},
      4/{Veszteségforgatókönyvek\\azonosítása}/{miért következik be az UCA,\\miért nem hajtódik végre\\a helyes akció}} {
    \node[l] (s\i) at ({(\i-1)*3.9},0) {\textbf{\t}\\[3pt]{\tiny \d}};
    \node[circle, fill=peKiemel, text=white, font=\scriptsize\bfseries, minimum size=5mm, inner sep=0] at ($(s\i.north)+(0,0.05)$) {\i};
  }
  \foreach \a/\b in {1/2,2/3,3/4} \draw[nyil] (s\a) -- (s\b);
  \draw[nyil, szaggatott] (s4.south) -- ++(0,-0.45) -| node[pos=0.25, below, font=\scriptsize] {követelmények, tervmódosítások, és szükség szerint a modell finomítása} (s1.south);
\end{tikzpicture}
\end{document}
